\section{Introduction}
\label{sec:introduction}
Large language model deployments often need outputs that are careful, factual, and safe. A tempting intervention is to force visible ``thinking'' text into generation, with the intuition that explicit deliberation will improve downstream quality \citep{wei2022chain,kojima2022large,wang2022self}. Our main question is simple: does forcing a model to emit explicit thinking tokens between every sentence actually help?

The practical stakes are high. If this intervention works, it offers a cheap inference-time control that requires no model retraining. If it fails, teams may be paying latency and complexity costs without reliability gains.

Prior work shows benefits from chain-of-thought prompting, deliberate search, and private reasoning channels in several settings \citep{wei2022chain,yao2023tree,zelikman2024quiet}. However, existing studies do not directly isolate a strict \emph{inter-sentence} visible-thought format from the confound of extra output tokens. We target this gap with an explicit token-budget control condition.

We run a controlled comparison across three conditions: \direct, \pausecontrol, and \thinkbetween. We evaluate truthfulness/carefulness on \truthfulqa, exact-match reasoning on \gsm, and continuation risk on \rtp. Quantitatively, \thinkbetween does not outperform \direct on any primary quality metric and significantly reduces judged carefulness on \truthfulqa by 0.40 points.

\para{{\bf what do we contribute?}} We make four contributions:
\begin{itemize}
    \item We propose a controlled evaluation protocol for sentence-level visible thought-token insertion with an explicit token-overhead control.
    \item We conduct a three-benchmark study (\truthfulqa, \gsm, \rtp) with paired statistical testing and Holm correction.
    \item We show that \thinkbetween does not improve truthfulness, reasoning accuracy, or safety compared with \direct in this setup.
    \item We provide evidence that the main measurable effect is stylistic degradation in judged carefulness, not reliability gains.
\end{itemize}

\para{{\bf how is the paper organized?}} \secref{sec:related} positions our study in prior literature, \secref{sec:method} details the protocol, \secref{sec:results} reports quantitative results, and \secref{sec:discussion} analyzes implications and limitations.
